\begin{abstract}
Fourier Neural Operators (FNOs) have achieved remarkable success in solving partial differential equations 
but lack built-in uncertainty quantification capabilities essential for safety-critical applications. 
This paper presents Evidential FNO, which integrates Evidential Deep Learning with FNOs to enable single-pass 
uncertainty estimation. We systematically compare three uncertainty quantification approaches 
(Monte Carlo Dropout, Deep Ensembles, and Evidential Deep Learning) across two benchmark problems: Navier-Stokes 
equations and Darcy flow. Our evaluation encompasses calibration quality, coverage guarantees, predictive accuracy, 
and uncertainty-error correlation. We investigate six regularization schemes for evidential learning and
identify that Inverse-Variance and Log-Barrier regularization achieve the best calibration with low Expected
Calibration Error and coverage near 93\%. Evidential methods provide these benefits with single-model,
single-pass inference, offering significant speedup over ensemble methods. We additionally assess whether the
predicted uncertainty is physically meaningful, by measuring its spatial alignment with the residual of the
governing equation, and find this alignment to be problem-dependent: on transient Navier-Stokes the most
uncertain regions coincide with the largest physics violations (rank correlation up to 0.63, with
$1.5$--$1.6\times$ residual concentration), whereas on the near-exactly-solved steady Darcy problem the
alignment essentially vanishes.
We stress-test this behavior on two harder regimes and find the same signature with opposite sign: on
shock-capturing compressible Euler flow the predicted uncertainty concentrates sharply on discontinuities
($3.9\times$ on the steepest-gradient cells), whereas under long-horizon autoregressive rollout it
\emph{collapses} even as error accumulates (anti-correlated with error at a mean of $-0.94$). Both trace to a
single mechanism: the pointwise evidence is a learned function of the current input's local structure,
calibrated only in-distribution, so it flags locally hard inputs but cannot sense global drift off the training
manifold. Our analysis also shows that out-of-distribution detection depends critically on the type of shift:
parameter extrapolation performs near random chance, while cross-equation detection achieves perfect
separation. These findings establish evidential methods as preferred for calibrated in-distribution prediction
while highlighting that ensemble methods remain necessary when robust epistemic uncertainty or out-of-distribution
detection is required. We further identify the pointwise independence assumption inherent to per-grid-point
evidential parameterization as the principal barrier to faithful uncertainty quantification over spatially
correlated physical fields.
\end{abstract}